% c11.fig.zerlegung — what Theorem c11.thm.separation says, as a diagram.
%
% A schematic, not a plot: no length in this picture is data.  It restates
%     U H_tot = H_cm + H_rel,   H_cm = -(hbar^2/2M) Lap_R,   H_rel = -(hbar^2/2mu) Lap_r + V(r)
% (DERIVATION.md Result 1.1): two particles of masses m_e, m_p coupled by V(r_e - r_p) are the
% same system as a free particle of mass M (coordinate R) together with one particle of mass mu
% (coordinate r) in the potential V(r).
% Greyscale: outlines and labels only.
\begin{tikzpicture}[x=1cm, y=1cm, >={Stealth[length=5pt,width=4pt]}]
  % left: the two particles
  \draw[tpc leader, rounded corners=3pt] (0,0) rectangle (3.5,3.2);
  \fill[tpc-plate-ink] (0.8,0.9) circle (3.4pt);
  \draw[tpc-plate-ink, fill=tpc-plate, line width=\tpcstrokeregular] (2.7,2.2) circle (2.4pt);
  \draw[tpc dynamics, dash pattern=on 3pt off 2pt] (0.95,1.0) -- (2.6,2.13);
  \node[tpc label, anchor=north] at (0.8,0.72) {$m_{\mathrm{p}}$, $\boldsymbol{r}_{\mathrm{p}}$};
  \node[tpc label, anchor=south] at (2.7,2.38) {$m_{\mathrm{e}}$, $\boldsymbol{r}_{\mathrm{e}}$};
  \node[tpc label, anchor=north west] at (1.75,1.5) {$V(\boldsymbol{r}_{\mathrm{e}} - \boldsymbol{r}_{\mathrm{p}})$};
  \node[tpc label, anchor=south] at (1.75,3.25) {$\op{H_{\mathrm{tot}}}$};
  % the change of coordinates
  \draw[tpc axis, ->] (3.75,1.6) -- (4.85,1.6) node[midway, above, tpc label] {$\op{U}$};
  % right: centre of mass
  \draw[tpc leader, rounded corners=3pt] (5.1,1.75) rectangle (10.1,3.2);
  \draw[tpc-plate-ink, line width=\tpcstrokeregular] (5.75,2.45) ++(-3pt,-3pt) -- ++(6pt,6pt)
    (5.75,2.45) ++(-3pt,3pt) -- ++(6pt,-6pt);
  \draw[tpc axis, ->] (5.95,2.45) -- (6.75,2.45);
  \node[tpc label, anchor=west, align=left] at (6.95,2.45)
    {$\op{H_{\mathrm{cm}}}$: \tpcifenglish{free particle,}{freies Teilchen,}\\
     \tpcifenglish{mass}{Masse} $M$, \tpcifenglish{coordinate}{Koordinate} $\boldsymbol{R}$};
  % right: relative motion
  \draw[tpc leader, rounded corners=3pt] (5.1,0) rectangle (10.1,1.45);
  \draw[tpc leader] (5.6,0.72) -- (5.9,0.72) (5.75,0.57) -- (5.75,0.87);
  \draw[tpc-plate-ink, fill=tpc-plate-muted, line width=\tpcstrokeregular] (6.6,1.0) circle (2.4pt);
  \draw[tpc dynamics, dash pattern=on 3pt off 2pt] (5.8,0.74) -- (6.52,0.97);
  \node[tpc label, anchor=west, align=left] at (6.95,0.72)
    {$\op{H_{\mathrm{rel}}}$: \tpcifenglish{one particle,}{ein Teilchen,}\\
     \tpcifenglish{mass}{Masse} $\mu$, \tpcifenglish{potential}{Potential} $V(\boldsymbol{r})$};
  \node[tpc label] at (7.6,1.6) {$+$};
\end{tikzpicture}
